% from ega4-2-fr.tex
\oldpage[IV-2]{184}
pour tout $i$; en vertu de (7.1.3.1), (7.1.3.2) et (7.1.3.4) et de
l'inégalité $\dim(Y'_i)\leqslant\dim(Y')$, on en déduit
\begin{equation}
\dim(X)\leqslant\dim(Y)+\operatorname{codim}(Y,X)
\tag{7.1.3.6}
\label{IV.2.eq.7.1.3.6-fr}
\end{equation}
donc les deux membres de cette inégalité sont égaux d'après (0, 14.2.2.2);
en outre cette égalité entraîne
\begin{equation}
\dim(Y'_i)=\dim(Y')=\dim(Y)
\tag{7.1.3.7}
\label{IV.2.eq.7.1.3.7-fr}
\end{equation}
pour tout $i$, ce qui prouve la seconde assertion de (ii).

Démontrons maintenant la première assertion de (ii). Soient $\mathfrak p_j$
($1\leqslant j\leqslant n$) les idéaux premiers de $A$ minimaux parmi ceux
qui contiennent $\mathfrak a$, et soient $\mathfrak p'_{jh}$ les idéaux
premiers de $A'$ minimaux parmi ceux contenant $\mathfrak p_jA'$
($1\leqslant h\leqslant m_j$); on sait alors (2.3.4) que les
$\mathfrak p'_{jh}$ ($1\leqslant j\leqslant n$, $1\leqslant h\leqslant
m_j$ pour tout $j$) sont aussi les idéaux premiers de $A'$ minimaux parmi
ceux qui contiennent $\mathfrak aA'$. D'après ce qu'on a vu plus haut, pour
tout $j$, les dimensions des anneaux $A'/\mathfrak p'_{jh}$
($1\leqslant h\leqslant m_j$) sont toutes égales et sont donc égales à
$\dim(A'/\mathfrak p_jA')$, elle-même égale à $\dim(A/\mathfrak p_j)$ par
(7.1.3.2). Dire que $A'/\mathfrak aA'$ est équidimensionnel signifie donc
que les dimensions des $A/\mathfrak p_j$ sont toutes égales, c'est-à-dire
que $A/\mathfrak a$ est équidimensionnel.

Prouvons enfin (i). Soit $\mathfrak r'$ un idéal premier de $A'$ minimal
parmi ceux contenant $\mathfrak mA'$, et soit $A''=A'_{\mathfrak r'}$;
comme $A''$ est un $A'$-module plat, c'est aussi un $A$-module plat; en
outre $A''$ est équidimensionnel et caténaire (0, 16.1.4), et
$\mathfrak mA''$ est un idéal de définition de $A''$. On peut donc se
borner au cas où $\mathfrak mA'$ est un idéal de définition de $A'$. Si
alors $\mathfrak p$ et $\mathfrak q\subset\mathfrak p$ sont deux idéaux
premiers de $A$, on peut appliquer (ii) à $A/\mathfrak q$ et à
$A'/\mathfrak qA'$, qui est équidimensionnel et plat sur
$A/\mathfrak q$; si $Z=\operatorname{Spec}(A/\mathfrak q)$ et
$Y=\operatorname{Spec}(A/\mathfrak p)$, on a donc
$\dim(Z)=\dim(Y)+\operatorname{codim}(Y,Z)$; en outre on peut appliquer
(7.1.3.6), qui montre que
$\operatorname{codim}(Y,X)=\dim X-\dim Y$; comme $X$ est
équicodimensionnel, ces relations montrent qu'il est biéquidimensionnel en
vertu de (0, 14.3.3).
\end{proof}

\begin{corollary}[7.1.4]
\label{IV.2.s7.1.4-fr}
Soit $A$ un anneau local noethérien formellement équidimensionnel. Alors:
\begin{enumerate}
\item[(i)] $A$ est équidimensionnel et caténaire (autrement dit,
biéquidimensionnel).
\item[(ii)] Soit $\mathfrak a$ un idéal de $A$; pour que $A/\mathfrak a$ soit
équidimensionnel, il faut et il suffit que $A/\mathfrak a$ soit formellement
équidimensionnel. En particulier, pour tout idéal premier $\mathfrak p$ de
$A$, $A/\mathfrak p$ est formellement équidimensionnel.
\end{enumerate}
\end{corollary}

\begin{proof}
Si $\mathfrak m$ est l'idéal maximal de $A$, $\mathfrak m\widehat A$ est
idéal de définition de $\widehat A$. Par hypothèse $\widehat A$ est
équidimensionnel, et on sait d'autre part qu'il est caténaire (5.6.4); il
suffit alors d'appliquer (7.1.3) à $A'=\widehat A$.
\end{proof}

\begin{corollary}[7.1.5]
\label{IV.2.s7.1.5-fr}
Soit $A$ un anneau local noethérien tel qu'il existe un $A$-module de type
fini $M$ qui soit un $A$-module de Cohen--Macaulay et pour lequel
$\operatorname{Supp}(M)=\operatorname{Spec}(A)$ (ce qui aura lieu par
exemple si $A$ est un anneau de Cohen--Macaulay). Alors $A$ est formellement
équidimensionnel, donc (7.1.4) pour qu'un anneau quotient $B$ de $A$ soit
formellement équidimensionnel, il faut et il suffit qu'il soit
équidimensionnel.
\end{corollary}

\begin{proof}
